\chapter{Metodología}

Este capítulo describe el método de modelado: la carga de los volúmenes y su
recorte de entrada, la arquitectura de la red convolucional 3D, la función de
pérdida y el esquema de optimización, el aumento de datos, la estrategia de
partición a nivel de sujeto y el entorno de cómputo. Todos los valores que se
indican proceden del código fuente o de los ficheros de configuración del
proyecto, y son comunes a los distintos regímenes de evaluación, de modo que las
diferencias observadas entre ellos no se deban a cambios en el método.

\section{Preprocesado homogéneo}

El capítulo anterior fijó el preprocesado común a todas las cohortes
(extracción cerebral, reorientación a RAS, remuestreo a 1\,mm\textsuperscript{3},
ajuste a \mbox{192\,$\times$\,224\,$\times$\,192} vóxeles y normalización
\emph{z-score} sobre el tejido), cuyo resultado es, para cada estudio, un fichero
\texttt{.npz} con los volúmenes T1 y T2 ya normalizados. La metodología de
modelado parte de ese punto.

En el momento de la carga, cada volumen T1 y T2 se apila en un tensor de dos
canales y se ajusta a la geometría de entrada de la red mediante un
\textbf{recorte centrado a 128\,$\times$\,160\,$\times$\,128} vóxeles. Conviene
distinguir este recorte del tamaño de preprocesado: el preprocesado deja todos
los volúmenes en una rejilla común de 192\,$\times$\,224\,$\times$\,192, y es en
la carga donde se recorta a 128\,$\times$\,160\,$\times$\,128 para reducir el
coste de memoria de las convoluciones 3D conservando la región cerebral central.
Si un volumen fuera menor que la geometría objetivo en algún eje, se rellena de
forma centrada con ceros. Este recorte se aplica de forma centrada tanto en
entrenamiento como en validación y test, lo que elimina cualquier discrepancia
sistemática entre particiones por el modo de recorte.

El calificativo de \emph{homogéneo} es deliberado: el mismo preprocesado y la
misma transformación de carga se aplican a las cuatro cohortes del pool y, sin
más cambio que el número de canales de entrada, a la cohorte intra-dominio. Así,
ninguna diferencia de resultados entre datasets puede atribuirse a un
tratamiento distinto de los datos.

\section{Arquitectura del modelo CNN 3D}

El modelo, denominado \texttt{BrainTumorCNN3D}, es una red convolucional 3D
compacta para clasificación binaria. Recibe un tensor de dos canales (T1 y T2) y
produce un único logit, cuya sigmoide se interpreta como probabilidad de la
clase tumor. Consta de un extractor de características con cuatro bloques
convolucionales y un clasificador denso.

Cada bloque convolucional (\texttt{ConvBlock3D}) encadena dos convoluciones 3D
de núcleo $3\times3\times3$ con relleno 1 y sin sesgo, cada una seguida de
normalización y activación ReLU, y un \emph{dropout} espacial al final del
bloque. El número de canales se duplica de bloque en bloque a partir de un ancho
base de 12: $12 \rightarrow 24 \rightarrow 48 \rightarrow 96$. Entre bloques
consecutivos se intercala un \texttt{MaxPool3d} de factor 2, que reduce la
resolución espacial a la mitad en cada eje; tras el último bloque, un
\texttt{AdaptiveAvgPool3d} colapsa la dimensión espacial a $1\times1\times1$ y
produce un vector de \textbf{96 dimensiones}, que constituye la representación
latente del estudio. El clasificador aplica sobre ese vector una capa lineal
$96 \rightarrow 96$ con ReLU y \emph{dropout}, y una capa lineal final
$96 \rightarrow 1$. La Tabla~\ref{tab:arquitectura} resume la estructura.

\begin{table}[ht]
  \centering
  \caption[Arquitectura de \texttt{BrainTumorCNN3D}]{Arquitectura de \texttt{BrainTumorCNN3D} con la geometría de entrada
  de 2\,$\times$\,128\,$\times$\,160\,$\times$\,128. Cada bloque convolucional
  contiene dos convoluciones $3\times3\times3$ con GroupNorm y ReLU; el
  \emph{dropout} indicado es el \texttt{Dropout3d} aplicado al final del bloque.}
  \label{tab:arquitectura}
  \begin{tabular}{l l c}
    \hline
    \textbf{Etapa} & \textbf{Operación} & \textbf{Canales / dim.} \\
    \hline
    Entrada            & Tensor T1+T2                       & 2 \\
    Bloque 1           & ConvBlock3D ($dropout=0{,}0$)      & 12 \\
                       & MaxPool3d (2)                      & 12 \\
    Bloque 2           & ConvBlock3D ($dropout=0{,}125$)    & 24 \\
                       & MaxPool3d (2)                      & 24 \\
    Bloque 3           & ConvBlock3D ($dropout=0{,}25$)     & 48 \\
                       & MaxPool3d (2)                      & 48 \\
    Bloque 4           & ConvBlock3D ($dropout=0{,}25$)     & 96 \\
    \emph{Pooling}     & AdaptiveAvgPool3d (1)              & 96 (embedding) \\
    \hline
    Clasificador       & Linear $96\!\rightarrow\!96$ + ReLU + Dropout(0,25) & 96 \\
                       & Linear $96\!\rightarrow\!1$        & 1 (logit) \\
    \hline
  \end{tabular}
\end{table}

Una decisión de diseño relevante es el uso de \textbf{normalización por grupos}
(GroupNorm, con 8 grupos) en lugar de la habitual normalización por lotes
(BatchNorm)~\cite{wu2018groupnorm}. El motivo es directo: con un tamaño de lote
de un único estudio, los estadísticos por lote de la BatchNorm son muy ruidosos
y sus estadísticos acumulados no convergen, lo que en la práctica genera una
brecha grande entre entrenamiento y validación; la GroupNorm normaliza por
grupos de canales dentro de cada muestra y es estable con lote unitario.

El modelo es deliberadamente pequeño: en la configuración de dos canales tiene
\textbf{504\,553} parámetros entrenables (504\,229 en la variante de un solo
canal). Esa misma arquitectura, sin cambios salvo en el número de canales de
entrada (un canal en lugar de dos), se reutiliza en el experimento intra-dominio
descrito en capítulos posteriores, de manera que la comparación entre regímenes
mantiene constante el modelo.

\section{Función de pérdida, optimización e hiperparámetros}

El problema se entrena como clasificación binaria con una sola salida. La función
de pérdida es la entropía cruzada binaria sobre el logit
(\texttt{BCEWithLogitsLoss}), que combina de forma numéricamente estable la
sigmoide y la entropía cruzada. Se emplea un peso de clase positiva
$\texttt{pos\_weight}=1{,}0$, es decir, sin reponderar las clases: el pool está
aproximadamente equilibrado y un peso unitario elimina la hipótesis de que un
sesgo de la pérdida favoreciera artificialmente la clase tumor.

La optimización se realiza con \textbf{AdamW}~\cite{loshchilov2019adamw}, con tasa
de aprendizaje $5\times10^{-5}$ y decaimiento de pesos $1\times10^{-3}$. La tasa
de aprendizaje sigue una planificación de tipo coseno
(\texttt{CosineAnnealingLR}) a lo largo de las épocas, lo que permite un ajuste
fino al final del entrenamiento. Los hiperparámetros completos se recogen en la
Tabla~\ref{tab:hiperparametros}.

\begin{table}[ht]
  \centering
  \caption[Hiperparámetros de entrenamiento]{Hiperparámetros de entrenamiento (tomados de \texttt{configs/train\_3d.yaml}
  y del bucle de entrenamiento).}
  \label{tab:hiperparametros}
  \begin{tabular}{l l}
    \hline
    \textbf{Hiperparámetro} & \textbf{Valor} \\
    \hline
    Optimizador                  & AdamW \\
    Tasa de aprendizaje          & $5\times10^{-5}$ \\
    Decaimiento de pesos         & $1\times10^{-3}$ \\
    Planificador                 & CosineAnnealingLR ($T_{\max}=40$) \\
    Épocas máximas               & 40 \\
    Tamaño de lote               & 1 \\
    Función de pérdida           & BCEWithLogitsLoss ($\texttt{pos\_weight}=1{,}0$) \\
    Precisión mixta (AMP)        & desactivada \\
    Paciencia (early stopping)   & 10 \\
    Criterio de selección        & balanced accuracy (val) con $\text{sen}\geq0{,}80$ \\
    Semilla                      & 42 \\
    \hline
  \end{tabular}
\end{table}

La \textbf{selección del modelo} merece detalle, por su relevancia metodológica.
En cada época se evalúa la validación y se calcula la \emph{balanced accuracy}
como media de sensibilidad y especificidad. Solo se guarda como mejor modelo el
checkpoint que mejora esa métrica \emph{y} mantiene una sensibilidad de
validación de al menos $0{,}80$; esta restricción evita conservar modelos
triviales (por ejemplo, uno que clasifique todo como sano, con especificidad
alta pero sensibilidad nula). El entrenamiento se detiene anticipadamente si
transcurren 10 épocas sin mejora bajo ese criterio. El \textbf{conjunto de test
se evalúa una única vez} al final, con el mejor checkpoint seleccionado en
validación, de modo que no interviene en la selección del modelo ni en el ajuste
de hiperparámetros.

\section{Aumento de datos}

El aumento de datos se aplica exclusivamente al conjunto de entrenamiento; la
validación y el test se procesan sin alteración. Se combinan transformaciones
geométricas y de intensidad:

\begin{itemize}
  \item \textbf{Reflexiones espaciales} en cada uno de los tres ejes, con
  probabilidad $0{,}5$ por eje, aplicadas de forma compartida a los canales T1 y
  T2 para preservar su correspondencia.
  \item \textbf{Corrección gamma} con factor muestreado uniformemente en
  $[0{,}8;\ 1{,}25]$ (probabilidad $0{,}5$), preservando el signo de los valores
  \emph{z-score}.
  \item \textbf{Ruido gaussiano} de baja amplitud (desviación típica $0{,}03$,
  probabilidad $0{,}5$).
\end{itemize}

Las transformaciones de intensidad (gamma y ruido) se aplican \textbf{únicamente
sobre los vóxeles no nulos}, es decir, sobre la máscara de cerebro resultante de
la extracción craneal, y no sobre el fondo. Esta decisión es un intento
metodológico explícito de mitigar la firma de dominio: el objetivo es forzar al
modelo a apoyarse en el \emph{contraste relativo} del tejido y reducir su
dependencia del brillo absoluto característico de cada cohorte, evitando además
inyectar señal en el fondo (relevante en una de las cohortes, que arrastra un
borde de fondo no nulo). Como se discutirá más adelante, esta mitigación a nivel
de aumento de datos no basta para neutralizar un confound que está presente en
la propia composición del conjunto.

\section{Estrategia de partición a nivel de sujeto}

La partición en entrenamiento, validación y test se realiza a nivel de
\emph{sujeto}, no de fichero, mediante una doble salvaguarda. En primer lugar,
los estudios se agrupan por \texttt{subject\_id} antes de repartirse, de manera
que todos los ficheros de un mismo sujeto caen \emph{por construcción} en la
misma partición; esto previene la fuga de información que se produciría si un
sujeto con varias adquisiciones (por ejemplo, distintas sesiones de una misma
cohorte) apareciese simultáneamente en entrenamiento y en test. En segundo
lugar, el reparto se estratifica conjuntamente por \texttt{(dataset, label)},
para que ningún dataset quede sub- o sobre-representado en validación o test por
efecto del azar. Las proporciones son $70\%$ / $15\%$ / $15\%$ y el reparto es
reproducible con semilla $42$.

El proyecto incluye, además, una auditoría específica de la partición que
comprueba el solapamiento de sujetos entre conjuntos. Conviene subrayar el
alcance correcto de esta salvaguarda: la partición por sujeto descarta la fuga
de información \emph{por partición}, pero no neutraliza el confound
dominio$\leftrightarrow$clase, que es una propiedad de la composición del conjunto (cada cohorte
aporta una sola clase) y no del modo de repartir los datos. Por ello el split,
aun siendo correcto, no salva por sí solo la evaluación, lo que motiva los
diseños experimentales del capítulo siguiente. Para el régimen
\emph{leave-one-dataset-out}, esta partición aleatoria se sustituye por una
partición por dataset generada específicamente, como se detalla más adelante.

\section{Entorno computacional y herramientas}

El sistema se implementó en Python con PyTorch como marco de aprendizaje
profundo~\cite{paszke2019pytorch}. El entrenamiento se ejecutó sobre GPU AMD
mediante la pila ROCm. Las versiones principales del entorno, fijadas en el
fichero de dependencias del proyecto, son PyTorch~2.5.1 (compilación
ROCm~6.2), MONAI~1.5.0, HD-BET~2.0.1, NumPy~2.1.3, nibabel~5.4.2 y
scikit-learn~1.6.1. La conversión inicial de los estudios en formato DICOM a
NIfTI, previa al preprocesado, se realizó con dcm2niix~\cite{li2016dcm2niix}.

Para acelerar las convoluciones 3D con tamaño de entrada estable se habilitaron
el \emph{benchmark} de cuDNN, la aritmética TF32 y el formato de memoria
\emph{channels-last} 3D; la precisión mixta automática se mantuvo desactivada,
ya que con lote unitario y GroupNorm la representación en \texttt{float16}
amplificaba inestabilidades numéricas. La reproducibilidad se sostiene en una
semilla fija ($42$) para la generación de la partición y el entrenamiento; los
análisis posteriores que requieren remuestreo (intervalos de confianza por
\emph{bootstrap} y proyecciones) emplean su propia semilla.

Por último, el régimen \emph{leave-one-dataset-out} se ejecuta con una
configuración (\texttt{configs/train\_lodo.yaml}) idéntica a la del entrenamiento
principal salvo en dos detalles: no regenera la partición ---lee la partición por
dataset preparada aparte--- y escribe en su propio directorio de salida. Mantener
constantes el modelo y todos los hiperparámetros garantiza que la comparación
entre el régimen confundido y el cross-dataset sea estrictamente
\emph{apples-to-apples}, atribuible solo al cambio de partición y no a diferencias
de método.
